\begin{abstract}
A language-model agent works inside a \emph{harness}: its prompt, memory, tools and helper code.
Continual Harness lets an agent rewrite its own harness while it works, and it improves large
frontier models. Its authors suggest that weaker models fall below a ``capability floor'' and cannot
improve themselves this way. We test this directly with two small models that run on a laptop,
\texttt{qwen3.5:0.8b} and \texttt{gemma4:e2b}, on 20 computer tasks (coding, files, shell and data)
with automatic pass/fail checks. We re-implement Continual Harness for local models and add five
techniques, each aimed at a weakness of small models: a separate, larger refiner; a regression gate
that keeps an edit only if recent tasks do not get worse; a memory database with retrieval; skills
that are kept only if their code runs; and structured error messages. We find three things. First,
Continual Harness as-is makes both models significantly worse than a fixed harness (qwen: 32\% to
15\%; gemma: 92\% to 77\%). Second, no single technique brings the small model back to the
fixed-harness level, but all five together do (qwen: 35\%, significantly better than Continual
Harness as-is, $p = 0.023$), with a large but not significant gain on held-out tasks (27\% to
60\%). Third, a simple Reflexion-style memory makes the smallest model much worse (5\%). For the
stronger model, the regression gate rejects most edits, which removes most of the harm but gives no
gain over the fixed harness. Self-editing below the capability floor is harmful by default, and
the main value of our techniques is to limit that harm. We release our code and tasks.
\end{abstract}
